\ndetext{joint \`a~\Ref{X.3.3} et aux crit\`eres~\Ref{XII.2.4} et~\Ref{XII.3.4}, on obtient le th\'eor\`eme de Lefschetz relatif suivant. Soit $f:X\to S$ un morphisme projectif et plat de sch\'emas noeth\'eriens connexes et soit $D$ un diviseur de Cartier relatif effectif dans $X$ et relativement ample. Si, pour tout $s\in S$, la profondeur de $X_s$ en chaque point ferm\'e est $\geq 2$, alors $D$ est connexe et, pour tout ouvert $U$ de $X$ contenant $D$, la fl\`eche $i_U:\pi_1(D)\to\pi_1(U)$ est surjective. Si de plus, la profondeur de $X_s$ le long de chaque point ferm\'e de $D_s$ est $\geq 3$ et si les anneaux locaux de $X$ en ses points ferm\'es sont purs, par exemple d'intersection compl\`ete (\cf \Ref{X}~\Ref{X.3.4}), alors $i_X$ est un isomorphisme. \Cf Bost~J.-B., {\og Lefschetz theorem for Arithmetic Surfaces\fg}, \emph{Ann. Sci. \'Ec. Norm. Sup. (4)} \textbf{32} (1999), p\ptbl 241-312, th\'eor\`emes 1.1 et 2.1. Dans le cas o\`u $X$ est simplement une surface projective lisse et g\'eom\'etriquement connexe sur un corps, on~a toujours connexit\'e de $D$ et surjectivit\'e de $\pi_1(D)\to\pi_1(U)$ (o\`u $U$ ouvert contenant~$D$) pour $D$ seulement nef de carr\'e $>0$ (\cf \loccit, th\'eor\`eme 2.3 et aussi le th\'eor\`eme 2.4 pour des surfaces seulement normales et compl\`etes). Dans le cas d'une surface arithm\'etique (normale et quasi-projective) $X$ \sisi{au dessus}{au-dessus} d'un anneau d'entiers $\Oo_K$, Bost, am\'eliorant des r\'esultats de Ihara (Ihara~Y., {\og Horizontal divisors on arithmetic surfaces associated with Bely\u\i\ uniformizations\fg}, in \emph{The Grothendieck theory of dessins d'enfants (Luminy, 1993)}, London Math. Soc. Lect. Note Series, vol.~200, Cambridge Univ. Press, Cambridge, 1994, 245--254 ou \loccit, corollaire 7.2), a montr\'e que si un point $P\in X(\Oo_K)$, qui joue le r\^{o}le du diviseur $D$ dans la situation g\'eom\'etrique, v\'erifie certaines conditions de positivit\'e, alors la fl\`eche $\pi_1(X)\to\pi_1(\Spec \Oo_K)$ d\'eduite de la projection \'etait inversible d'inverse la fl\`eche $\pi_1(\Spec \Oo_K)\to\pi_1(X)$ d\'eduite de $P$ (\loccit, th\'eor\`eme 1.2).}